\chapter{Коллекции, ввод-вывод и многопоточность}
\label{ch:05}

Вы умеете описывать данные классами и защищать код исключениями. Но программа,
которая ничего не хранит, ничего не читает с диска и делает всё по очереди,
остаётся учебной. Глава добавляет три способности: хранить наборы данных в
\term{коллекциях}, работать с файлами через \term{потоки ввода-вывода} и делать
несколько дел сразу в \term{потоках выполнения}. Русское «поток» обозначает и
поток данных (\eng{stream}), и поток выполнения (\eng{thread}); первый называем
потоком ввода-вывода, второй~--- просто потоком.

\section{Коллекции вместо массивов}

Массив имеет фиксированный размер: объявив \texttt{new int[10]}, вы навсегда
получили десять ячеек. Это книжная полка на десять книг: чтобы поставить
одиннадцатую, придётся собрать полку побольше и переставить всё. Рутину берёт на
себя \term{каркас коллекций} (\eng{Java Collections Framework}).

Продолжим аналогию: \texttt{List}~--- резиновая полка, которая растёт и помнит
порядок; \texttt{Set}~--- стопка уникальных книг; \texttt{Queue}~--- очередь в
буфете; \texttt{Map}~--- каталог, по шифру находящий книгу
(рис.~\ref{fig:05-1}). В корне стоит \texttt{Iterable}, дающий общую для всех
коллекций способность обойти элементы циклом \texttt{for-each}.

\begin{figure}[htbp]\centering
\begin{forest} hierarchy
[Iterable
  [Collection
    [List [ArrayList] [LinkedList]]
    [Set [HashSet] [LinkedHashSet] [SortedSet [TreeSet]]]
    [Queue [PriorityQueue] [BlockingQueue] [Deque [ArrayDeque]]]
  ]
]
\end{forest}
\caption{Иерархия интерфейсов и реализаций коллекций}\label{fig:05-1}
\end{figure}

\texttt{LinkedList} хранит элементы цепочкой узлов: вставка в середину дёшева,
а доступ по индексу медленный; он реализует ещё и \texttt{Deque}~--- очередь с
доступом к обоим концам. \texttt{BlockingQueue} умеет задерживать поток, пока в
ней не появится место или элемент. Отдельной ветвью стоит \texttt{Map}
(\texttt{HashMap}, \texttt{LinkedHashMap}, \texttt{TreeMap}): он \emph{не}
наследует \texttt{Collection}, так как хранит не элементы, а пары «ключ~---
значение».

Имена реализаций читаются без справочника: сначала способ хранения, потом
интерфейс (\texttt{ArrayList}, \texttt{HashSet}), а приставка \texttt{Linked}
(\texttt{LinkedHashSet}, \texttt{LinkedHashMap}) добавляет к хеш-структуре
память о порядке вставки.

\section{Выбор реализации}

От выбора реализации зависят и скорость программы, и порядок, в котором придут
элементы (табл.~\ref{tab:05-1}). Особенно важен второй столбец: у
\texttt{HashSet} и \texttt{HashMap} порядок не определён вовсе и может
измениться при добавлении элементов. Ни одна реализация не потокобезопасна.

\begin{table}[htbp]
\caption{Основные реализации коллекций}
\label{tab:05-1}
\centering
\begin{tabular}{L{2.7cm}L{3.2cm}C{1.7cm}L{6.6cm}}
\toprule
Класс & Порядок элементов & Дубли\-каты & Когда применять \\
\midrule
\texttt{ArrayList} & по индексу & да & частое чтение по индексу \\
\texttt{HashSet} & не определён & нет & быстрая проверка наличия \\
\texttt{TreeSet} & по значению & нет & уникальность плюс сортировка \\
\texttt{HashMap} & не определён & нет & быстрый поиск по ключу \\
\texttt{TreeMap} & по ключу & нет & отсортированные ключи \\
\bottomrule
\end{tabular}
\end{table}

В листинге~\ref{lst:05-1} одни и те же данные ведут себя по-разному в
зависимости от того, куда их положили: обычная \texttt{Queue} отдаёт элементы
методом \texttt{poll()} в порядке прихода, а \texttt{PriorityQueue}~---
наименьший первым.

\begin{listing}[H]
\begin{minted}{java}
List<String> langs = new ArrayList<>(List.of("Java", "Python"));
langs.add("Java");                       // список допускает повтор
System.out.println(langs.size());        // 3

Set<String> unique = new HashSet<>(langs);
System.out.println(unique.size());       // 2: повтор отброшен

Map<String, Integer> stock = new HashMap<>();
stock.put("apples", 3);
stock.put("apples", 5);                  // перезапись, не пара
System.out.println(stock.get("apples")); // 5

// PriorityQueue - минимальная куча: порядок прихода не важен
PriorityQueue<Integer> pq = new PriorityQueue<>();
pq.offer(5); pq.offer(1); pq.offer(3);
System.out.print(pq.poll());             // 1, наименьший
\end{minted}
\caption{Список, множество, отображение и очередь}
\label{lst:05-1}
\end{listing}

\begin{oshibka}
\textbf{Частая ошибка: свой класс в \texttt{HashSet} или ключом
\texttt{HashMap}.} Хеш-коллекции ищут элемент по значению \texttt{hashCode()}, а
найденное сверяют методом \texttt{equals()}. Переопределив только
\texttt{equals()}, вы получите два «равных» объекта в разных корзинах, и
множество примет обе копии. Эти два метода переопределяют вместе и по одним и
тем же полям.
\end{oshibka}

\section{Порядок элементов: Comparable и Comparator}

Чтобы отсортировать коллекцию, надо объяснить Java, какой из двух объектов
меньше. \texttt{Comparable} встраивает правило внутрь класса: объект «знает», как
сравнить себя с себе подобным, и такой порядок называют \term{естественным}.
\texttt{Comparator}~--- стратегия сравнения снаружи, и таких стратегий может быть
сколько угодно. В листинге~\ref{lst:05-2} и дальше группа~--- это
\texttt{List<Student> group}, а студент описан записью: её компоненты сразу дают
методы \texttt{name()}, \texttt{faculty()} и \texttt{grade()}.

\begin{listing}[H]
\begin{minted}{java}
record Student(String name, String faculty, int grade)
        implements Comparable<Student> {

    // Естественный порядок: по возрастанию оценки
    @Override
    public int compareTo(Student other) {
        return Integer.compare(grade, other.grade);
    }
}

List<Student> group = new ArrayList<>(/* ... */);
Collections.sort(group);     // использует compareTo

// По имени, а при совпадении имён - по оценке
group.sort(Comparator.comparing(Student::name)
        .thenComparingInt(Student::grade));
\end{minted}
\caption{Естественный порядок и внешняя стратегия сравнения}
\label{lst:05-2}
\end{listing}

Метод \texttt{compareTo} возвращает отрицательное число, ноль или положительное.
Писать \texttt{this.grade - other.grade} не стоит: на больших числах разность
переполняет \texttt{int} и знак получается обратным.

\section{Обработка коллекций: Stream API}

Цикл описывает, \emph{как} обойти коллекцию. Появившийся в Java~8
\term{Stream API} позволяет описать, \emph{что} нужно получить: цепочка состоит
из источника, любого числа \term{промежуточных} операций (каждая возвращает новый
поток и ничего не вычисляет) и одной \term{терминальной}, запускающей обработку.
Готовые сборщики для \texttt{collect()} лежат в классе \texttt{Collectors}
пакета \texttt{java.util.stream} (листинг~\ref{lst:05-3}).

\begin{listing}[H]
\begin{minted}{java}
List<Integer> numbers = List.of(2, 3, 6, 7, 8, 10);

List<Integer> evenSquares = numbers.stream()
        .filter(n -> n % 2 == 0)   // промежуточная
        .map(n -> n * n)           // промежуточная
        .toList();                 // терминальная
System.out.println(evenSquares);   // [4, 36, 64, 100]

// Средний балл по каждому факультету
Map<String, Double> avg = group.stream()
        .collect(Collectors.groupingBy(
                Student::faculty,
                Collectors.averagingInt(Student::grade)));
\end{minted}
\caption{Фильтрация, преобразование и группировка}
\label{lst:05-3}
\end{listing}

Поток одноразовый: после терминальной операции нужен новый вызов
\texttt{stream()}. Исходную коллекцию он не меняет; правка списка во время
обработки даст \texttt{ConcurrentModificationException}. А
\texttt{parallelStream()} раскидывает работу по ядрам, но на десятке элементов
накладные расходы съедят выгоду.

\section{Потоки ввода-вывода}

Любой обмен данными с внешним миром~--- файлом, консолью, сетью~--- устроен
одинаково: данные идут потоком, порция за порцией, как вода по трубе.

\subsection{Байтовые и символьные потоки}

Потоки делятся на две ветви. \term{Байтовые}~--- наследники
\texttt{InputStream} и \texttt{OutputStream}~--- переносят сырые байты и годятся
для любых данных: изображений, архивов. \term{Символьные}~--- наследники
\texttt{Reader} и \texttt{Writer}~--- переносят символы и сами перекодируют их,
поэтому текст читают ими (рис.~\ref{fig:05-3}).

\begin{figure}[htbp]\centering
\begin{forest} hierarchy
[java.io
  [InputStream [FileInputStream] [ObjectInputStream]]
  [OutputStream [FileOutputStream] [PrintStream]]
  [Reader [FileReader] [BufferedReader]]
  [Writer [FileWriter] [BufferedWriter]]
]
\end{forest}
\caption{Основные классы потоков ввода-вывода}\label{fig:05-3}
\end{figure}

Схема подсказывает главный приём~--- \emph{обёртывание}. Классы с приставкой
\texttt{File} ходят на диск напрямую, по байту или символу, что очень медленно;
\texttt{Buffered}-классы оборачивают любой поток и дёргают диск раз на несколько
килобайт, а \texttt{BufferedReader} вдобавок приносит \texttt{readLine()}.
Кстати, \texttt{System.out}~--- это \texttt{PrintStream}, а
\texttt{System.in}~--- \texttt{InputStream}.

Открытый поток обязательно закрывают: иначе останется занят файловый дескриптор,
а буфер не попадёт на диск. Конструкция \texttt{try-with-resources} закрывает
ресурсы сама (листинг~\ref{lst:05-4}). В байтовой половине запомните идиому
цикла: \texttt{read(buf)} возвращает число реально прочитанных байт, а
\texttt{-1} означает конец файла, поэтому записывают ровно \texttt{n} байт.

\begin{listing}[H]
\begin{minted}{java}
// Байтовый поток: копируем файл порциями по 4 Кбайт
try (FileInputStream in = new FileInputStream("photo.jpg");
     FileOutputStream out = new FileOutputStream("copy.jpg")) {
    byte[] buf = new byte[4096];
    int n;
    while ((n = in.read(buf)) != -1) {
        out.write(buf, 0, n);
    }
}

// Символьный поток: пишем и читаем текст
try (BufferedWriter bw = new BufferedWriter(
        new FileWriter("text.txt", StandardCharsets.UTF_8))) {
    bw.write("Первая строка");
    bw.newLine();
}   // close() вызовется сам и сбросит буфер на диск

try (BufferedReader br = new BufferedReader(
        new FileReader("text.txt", StandardCharsets.UTF_8))) {
    String line;
    while ((line = br.readLine()) != null) {
        System.out.println(line);
    }
}
\end{minted}
\caption{Побайтовое копирование и построчное чтение текста}
\label{lst:05-4}
\end{listing}

Кодировка задана не случайно: без второго аргумента \texttt{FileWriter} и
\texttt{FileReader} берут кодировку платформы~--- на русской Windows это cp1251,
и файл, открытый UTF-8-редактором, покажет кракозябры. Константы кодировок~---
в классе \texttt{StandardCharsets} пакета \texttt{java.nio.charset}.

\subsection{Сериализация объектов}

Как сохранить в файл целый объект со всеми полями? Для этого есть
\term{сериализация} (\eng{serialization})~--- превращение объекта в
последовательность байтов и обратно; класс должен реализовать интерфейс-маркер
\texttt{Serializable} (листинг~\ref{lst:05-5}).

\begin{listing}[H]
\begin{minted}{java}
class Person implements Serializable {
    private static final long serialVersionUID = 1L;
    private String name;
    private transient String password;  // не сохраняется

    Person(String n, String p) { name = n; password = p; }
}

try (ObjectOutputStream oos = new ObjectOutputStream(
        new FileOutputStream("person.dat"))) {
    oos.writeObject(new Person("Ivan", "secret"));
}

try (ObjectInputStream ois = new ObjectInputStream(
        new FileInputStream("person.dat"))) {
    Person p = (Person) ois.readObject();
    // поле password восстановится как null
}
\end{minted}
\caption{Сохранение объекта в файл и чтение обратно}
\label{lst:05-5}
\end{listing}

Поле \texttt{serialVersionUID}~--- номер версии формата: изменится класс, а номер
нет~--- и чтение старого файла даст \texttt{InvalidClassException}.
\texttt{transient} помечает поле, которое сохранять не нужно; после чтения оно
получит \texttt{null} или ноль. Метод \texttt{readObject()} возвращает
\texttt{Object} (результат приводят к типу) и бросает проверяемое
\texttt{ClassNotFoundException}: метод чтения объявляют \texttt{throws
IOException, ClassNotFoundException}.

\subsection{Файловая система: NIO.2}

Операции с самими файлами~--- копирование, обход каталогов~--- на потоках
громоздки. В Java~7 появился пакет \texttt{java.nio.file}: путь~--- это
\texttt{Path}, а действия собраны в классе \texttt{Files}
(листинг~\ref{lst:05-6}).

\begin{listing}[H]
\begin{minted}{java}
Path path = Path.of("data", "notes.txt");
Files.createDirectories(Path.of("data"));
Files.writeString(path, "первая строка");
System.out.println(Files.readString(path) + " " + Files.size(path));

// Рекурсивный обход дерева каталогов
try (Stream<Path> tree = Files.walk(Path.of("data"))) {
    tree.filter(p -> p.toString().endsWith(".txt"))
        .forEach(System.out::println);
}
\end{minted}
\caption{Операции с файлами и каталогами через NIO.2}
\label{lst:05-6}
\end{listing}

Метод \texttt{createDirectories} создаёт и целевой каталог, и недостающие
промежуточные, а \texttt{createDirectory} требует существующего родителя. Метод
\texttt{Files.walk} обходит дерево рекурсивно и возвращает
\texttt{Stream<Path>}~--- вот и встреча двух тем; обходимый каталог должен
существовать, иначе будет \texttt{NoSuchFileException}.

\section{Многопоточность}

Один повар готовит обед по очереди: режет овощи, варит суп, жарит мясо. Два
повара делают то же параллельно, и обед готов раньше. Так и потоки делят работу,
загружая все ядра и не давая интерфейсу замереть на долгой операции.

\subsection{Процессы и потоки}

\term{Процесс} (\eng{process})~--- запущенная программа с собственным участком
памяти: команда \texttt{java MyApp} создаёт процесс с отдельной виртуальной
машиной. \term{Поток} (\eng{thread})~--- лёгкая единица работы \emph{внутри}
процесса: потоки делят общую кучу, но у каждого свой стек вызовов
(рис.~\ref{fig:05-4}). Общая память~--- источник всех бед: если два потока
правят одну переменную, результат непредсказуем.

\begin{figure}[htbp]\centering
\begin{tikzpicture}[font=\small]
\node[boxw=24mm] (t1) {Поток \texttt{main}\\[1pt]\scriptsize свой стек};
\node[boxw=24mm,right=4mm of t1] (t2) {Поток 2\\[1pt]\scriptsize свой стек};
\node[boxw=24mm,right=4mm of t2] (t3) {Поток 3\\[1pt]\scriptsize свой стек};
\node[boxw=84mm,below=10mm of t2] (hp)
  {Общая куча: объекты, статические поля, массивы};
\draw[-Latex,dashed] (t1)--(hp);
\draw[-Latex,dashed] (t2)--(hp);
\draw[-Latex,dashed] (t3)--(hp);
\node[draw,rounded corners=3pt,fit=(t1)(t3)(hp),inner sep=3mm] {};
\end{tikzpicture}
\caption{Процесс: общая куча и собственный стек у каждого потока}\label{fig:05-4}
\end{figure}

\subsection{Создание потоков}

Поток создают двумя способами: наследовать \texttt{Thread} и переопределить
\texttt{run()} либо реализовать \texttt{Runnable} и передать задачу в конструктор
\texttt{Thread} (листинг~\ref{lst:05-7}). Второй предпочтителен: наследник
\texttt{Thread} израсходовал единственное наследование и смешал задание с
механизмом его выполнения.

\begin{listing}[H]
\begin{minted}{java}
// Способ 1: наследование Thread
class Loader extends Thread {
    @Override public void run() { System.out.println(getName()); }
}
new Loader().start();

// Способ 2: реализация Runnable, обычно лямбдой
Runnable task = () -> System.out.println("задача пошла");
Thread worker = new Thread(task, "worker");
worker.start();   // именно start(), а не run()
worker.join();    // главный поток ждёт worker
\end{minted}
\caption{Два способа создать поток выполнения}
\label{lst:05-7}
\end{listing}

\begin{oshibka}
\textbf{Частая ошибка: вызов \texttt{run()} вместо \texttt{start()}.} Метод
\texttt{start()} создаёт новый поток и уже в нём вызывает \texttt{run()}. Прямой
вызов \texttt{run()}~--- обычный вызов метода: код выполнится в текущем потоке,
последовательно, и никакого параллелизма не будет. Программа при этом работает и
ошибок не выдаёт~--- тем и опасна.
\end{oshibka}

Метод \texttt{join()} заставляет вызывающий поток дождаться чужого завершения;
без него \texttt{main} закончится раньше, чем рабочие потоки допечатают вывод.
\texttt{Thread.sleep(мс)} усыпляет \emph{текущий} поток. Ещё два свойства задают
до \texttt{start()}: \texttt{setDaemon(true)} делает поток \term{фоновым}
(\eng{daemon})~--- машина не ждёт его и закрывается, как только закончились все
обычные потоки; \texttt{setPriority(int)} принимает число от~1 до~10, но это
лишь подсказка планировщику операционной системы, а не гарантия~--- если порядок
важен для правильности, нужна синхронизация.

\subsection{Жизненный цикл потока}

От создания до завершения поток проходит несколько состояний; увидеть текущее
позволяет метод \texttt{getState()} (рис.~\ref{fig:05-5}).

\begin{figure}[htbp]\centering
\begin{tikzpicture}[font=\small,
  st/.style={draw,rounded corners=7pt,font=\ttfamily\footnotesize,
    inner sep=4pt,minimum height=7mm,minimum width=26mm,align=center},
  lb/.style={font=\scriptsize,align=center}]
\node[st] (nw) {NEW};
\node[st,below=13mm of nw] (rn) {RUNNABLE};
\node[st,below=15mm of rn] (tm) {TERMINATED};
\node[st,right=34mm of rn,text width=30mm] (wt)
  {BLOCKED\\WAITING\\TIMED\_WAITING};
\draw[flowarrow] (nw)--(rn) node[midway,right,lb] {~start()};
\draw[flowarrow] (rn)--(tm) node[midway,right,lb] {~run() завершён};
\draw[-Latex] ([yshift=3mm]rn.east)--([yshift=3mm]wt.west)
  node[midway,above,lb] {wait(), sleep(),\\занятый монитор};
\draw[-Latex] ([yshift=-3mm]wt.west)--([yshift=-3mm]rn.east)
  node[midway,below,lb] {notify(), таймаут,\\монитор освобождён};
\end{tikzpicture}
\caption{Состояния потока выполнения}\label{fig:05-5}
\end{figure}

Созданный объект \texttt{Thread} находится в \texttt{NEW}, после
\texttt{start()} переходит в \texttt{RUNNABLE}: выполняется либо ждёт очереди на
процессоре. \texttt{BLOCKED}~--- ожидание занятого монитора,
\texttt{WAITING}~--- бессрочное после \texttt{wait()} или \texttt{join()},
\texttt{TIMED\_WAITING}~--- с таймаутом. Из любого ожидания поток возвращается в
\texttt{RUNNABLE}, а не завершается напрямую; в \texttt{TERMINATED} он попадает,
когда \texttt{run()} закончился.

Остановить чужой поток силой нельзя: оборванный на середине, он оставил бы данные
в противоречивом состоянии. Применяют \term{прерывание}: \texttt{interrupt()}
поднимает флаг, а поток сам решает, когда завершиться, поэтому долгую работу
пишут циклом \texttt{while (!Thread.currentThread().isInterrupted()) \{ ... \}}.
При перехвате \texttt{InterruptedException} флаг сбрасывается, и в \texttt{catch}
его восстанавливают вызовом \texttt{Thread.currentThread().interrupt()}.

\subsection{Состояние гонки и synchronized}

Самая коварная проблема многопоточности~--- \term{состояние гонки}
(\eng{race condition}). Строка \texttt{count++} состоит из трёх действий:
прочитать, прибавить, записать. Если между чтением и записью вмешается второй
поток, оба прочитают одно число и запишут один результат~--- инкремент пропадёт
(листинг~\ref{lst:05-8}).

\begin{listing}[H]
\begin{minted}{java}
class Counter {
    private int count = 0;

    // Гонка: три действия, между которыми
    // может вклиниться другой поток
    public void increment() { count++; }

    // Безопасно: монитор пропускает внутрь
    // только один поток одновременно
    public synchronized void safeIncrement() { count++; }

    public synchronized int get() { return count; }
}
\end{minted}
\caption{Небезопасный и безопасный счётчик}
\label{lst:05-8}
\end{listing}

Запустите десять потоков по тысяче инкрементов и вместо десяти тысяч получите
меньше, причём каждый запуск даст своё число. Иногда результат верен~--- потому
такие ошибки и тяжело ловить. Лекарство~--- \texttt{synchronized}: в каждый
объект Java встроен \term{монитор}, невидимый замок, и вошедший в
синхронизированный метод поток захватывает его, а остальные ждут. Гарантий
две~--- \term{взаимное исключение} и \term{видимость памяти}: записанное до
выхода из блока увидит следующий вошедший.

\subsection{volatile и атомарные классы}

Иногда нужна только видимость, без блокировки. Тогда поле помечают словом
\texttt{volatile}: чтение берёт значение из основной памяти, запись немедленно
туда попадает. Без пометки поток вправе держать значение в регистре и не
заметить чужого изменения (листинг~\ref{lst:05-9}).

\begin{listing}[H]
\begin{minted}{java}
class Worker {
    // без volatile поток не увидит, что флаг сброшен
    private volatile boolean running = true;

    public void stop() { running = false; }

    public void work() { while (running) { /* работа */ } }
}

// Счётчик без блокировок: аппаратная операция CAS
AtomicInteger counter = new AtomicInteger();
counter.incrementAndGet();      // атомарно
\end{minted}
\caption{Флаг остановки и атомарный счётчик}
\label{lst:05-9}
\end{listing}

Атомарности \texttt{volatile} не даёт: \texttt{count++} остаётся тремя
действиями. Для счётчиков берут атомарные классы (\texttt{AtomicInteger},
\texttt{AtomicLong}) из пакета \texttt{java.util.concurrent.atomic}~--- учтите,
что \texttt{import java.util.concurrent.*} этот подпакет не подхватывает.
Работают они на аппаратной инструкции \term{CAS} (\eng{compare-and-swap}): она
меняет значение, только если оно не изменилось с момента чтения, иначе повторяет
попытку. Правило выбора простое: \texttt{volatile}~--- для флагов, атомарные
классы~--- для счётчиков, \texttt{synchronized}~--- когда операцию из нескольких
действий надо выполнить целиком.

\subsection{Взаимодействие потоков: wait и notify}

Взаимного исключения мало, когда потокам нужно договариваться: один готовит
данные, другой забирает. В классе \texttt{Object}~--- заметьте, не в
\texttt{Thread}~--- определены \texttt{wait()}, \texttt{notify()} и
\texttt{notifyAll()}: первый освобождает монитор и усыпляет поток, второй будит
один ожидающий, третий~--- все. Вызывать их можно только внутри
синхронизированного блока, иначе получите \texttt{IllegalMonitorStateException}
(листинг~\ref{lst:05-10}).

\begin{listing}[H]
\begin{minted}{java}
class Buffer {
    private final Queue<Integer> items = new LinkedList<>();
    private final int capacity = 5;

    public synchronized void put(int x)
            throws InterruptedException {
        while (items.size() == capacity) {
            wait();          // полон - ждём потребителя
        }
        items.offer(x);
        notifyAll();         // будим потребителей
    }
    // take() симметричен: ждёт, пока items не опустеет,
    // снимает элемент и вызывает notifyAll()
}
\end{minted}
\caption{Буфер производителя и потребителя}
\label{lst:05-10}
\end{listing}

Обратите внимание на \texttt{while}, а не \texttt{if}, перед \texttt{wait()}:
машина вправе разбудить поток и без \texttt{notify()} (ложное пробуждение), да и
пока разбуженный заново захватывает монитор, другой успевает опустошить буфер.
Сравните с \texttt{sleep()}: тот усыпляет поток, но монитор \emph{не}
отпускает~--- спать внутри синхронизированного блока значит держать замок и
никого не пускать. Вручную эту машинерию не пишут: у \texttt{BlockingQueue}
методы \texttt{put()} и \texttt{take()} блокируют сами.

\subsection{Пулы потоков}

Создавать \texttt{new Thread()} на каждую задачу расточительно: поток занимает
от полумегабайта под стек, а при тысячах задач система исчерпает ресурсы.
Решение~--- \term{пул потоков}: готовый набор рабочих потоков разбирает задачи
из общей очереди (листинг~\ref{lst:05-11}).

\begin{listing}[H]
\begin{minted}{java}
ExecutorService pool = Executors.newFixedThreadPool(4);

// Задача без результата - Runnable
pool.execute(() -> System.out.println("готово"));

// Задача с результатом - Callable
Future<String> future = pool.submit(() -> {
    Thread.sleep(1000);
    return "результат вычисления";
});
System.out.println(future.get());  // ждёт завершения

pool.shutdown();  // обязательно: иначе JVM не завершится
\end{minted}
\caption{Пул потоков и получение результата задачи}
\label{lst:05-11}
\end{listing}

Фабрика \texttt{Executors} даёт и другие пулы: растущий по надобности
\texttt{newCachedThreadPool()} и однопоточный
\texttt{newSingleThreadExecutor()}. Объект \texttt{Future}~--- «обещание»
результата: \texttt{get()} блокирует до готовности, \texttt{cancel()} отменяет.

\begin{vrezka}
\textbf{Важно.} Забыв вызвать \texttt{shutdown()}, вы получите программу, которая
напечатала весь вывод, но не завершается: потоки пула не фоновые, и виртуальная
машина честно их ждёт.
\end{vrezka}

\subsection{Потокобезопасные коллекции}

Теперь понятно, почему реализации из табл.~\ref{tab:05-1} небезопасны:
\texttt{ArrayList} при одновременной записи из двух потоков теряет элемент или
падает посреди перестройки внутреннего массива. Обёртка
\texttt{Collections.synchronizedList(...)} заворачивает каждый метод в
\texttt{synchronized}~--- просто, но медленно, и защищает лишь отдельные вызовы,
а не обход: цикл \texttt{for-each} по такой обёртке надо самому заключить в блок
\texttt{synchronized (syncList)}, иначе получите
\texttt{ConcurrentModificationException}. Лучше взять
\texttt{java.util.concurrent}: \texttt{ConcurrentHashMap} блокирует не всю карту,
а отдельные участки, а \texttt{CopyOnWriteArrayList} копирует массив при записи,
зато чтение обходится без блокировок.

\praktikum{}

Задания выполняются на JDK~17 или новее; дополнительные библиотеки не нужны.

\subsection*{Задание 1. Коллекции: задачи, контакты, текст}

Класс \texttt{TaskManager} поверх \texttt{ArrayList<String>}: добавление в конец
и на позицию (при недопустимой~--- в конец), удаление по названию, получение по
индексу или \texttt{null} вне диапазона, сортировка по алфавиту.

Класс \texttt{PhoneBook} на \texttt{TreeMap<String, List<String>>}: добавление
номера к списку по имени, поиск по части имени без учёта регистра, удаление
номера (без номеров контакт исчезает целиком). Объясните, что изменится при
замене \texttt{TreeMap} на \texttt{HashMap}.

Класс \texttt{TextAnalyzer}: \texttt{getUniqueWords} возвращает слова текста в
нижнем регистре как \texttt{TreeSet<String>} (он же их отсортирует),
\texttt{commonWords}~--- пересечение слов двух текстов,
\texttt{wordFrequency}~--- \texttt{Map<String, Integer>} частот. Разбивайте
вызовом \texttt{split("[\textbackslash\textbackslash s.,!?;:()-]+")}: обратный
слеш в строковом литерале Java удваивается.

\subsection*{Задание 2. Аналитика студентов}

Наберите данные из листинга~\ref{lst:05-12} и напишите шесть запросов, каждый
одной цепочкой Stream API, без ручных циклов.

\begin{listing}[H]
\begin{minted}{java}
record Student(String name, String faculty,
               int grade, int year) {}

List<Student> students = List.of(
    new Student("Анна Иванова", "ИТ", 95, 2),
    new Student("Борис Петров", "Математика", 78, 3),
    new Student("Вера Сидорова", "ИТ", 88, 1),
    new Student("Егор Фёдоров", "Физика", 85, 3));
\end{minted}
\caption{Исходные данные для задания 2}
\label{lst:05-12}
\end{listing}

Требуются: студенты «ИТ» с оценкой не ниже 85, отсортированные по имени; средний
балл по факультетам (\texttt{groupingBy} с \texttt{averagingInt}); число
студентов на курсе (\texttt{counting}); лучший студент факультета
(\texttt{maxBy}); имена в одной строке через запятую (\texttt{joining});
статистика оценок (\texttt{summaryStatistics()}). Ответьте письменно, чем
\texttt{groupingBy} лучше ручной группировки циклом.

\subsection*{Задание 3. Файлы и сериализация}

Класс \texttt{FileOperations} со статическими методами: запись списка строк
через \texttt{PrintWriter} поверх \texttt{BufferedWriter}; чтение строк обратно
через \texttt{BufferedReader}; подсчёт строк, слов и символов; побайтовое
копирование файла с буфером. Всюду применяйте \texttt{try-with-resources},
методы объявляйте \texttt{throws IOException}.

Вторая часть~--- сериализация. Класс \texttt{StudentRecord} с полями имени,
оценки, факультета и пароля: пароль пометьте \texttt{transient}, задайте
\texttt{serialVersionUID}. Сохраните два объекта в файл и прочитайте обратно,
объявив метод чтения \texttt{throws IOException, ClassNotFoundException}.
Объясните, почему пароль оказался равен \texttt{null} и что будет, если добавить
в класс поле, не меняя \texttt{serialVersionUID}.

\subsection*{Задание 4. Потоки выполнения и синхронизация}

Смоделируйте загрузку файлов: три потока «скачивают» файлы разного размера,
каждые 100~мс печатая свой прогресс. Объясните письменно, почему строки разных
потоков перемешиваются и что изменится, если убрать вызовы \texttt{join()}.

Реализуйте оба счётчика из листинга~\ref{lst:05-8} и метод на десять потоков по
тысяче инкрементов; прогоните тест по пять раз для каждого счётчика. Всегда ли
ошибается небезопасный вариант?

Соберите пару «производитель~--- потребитель» на \texttt{ArrayBlockingQueue}
ёмкостью~5: производитель кладёт десять задач с паузой 100~мс, потребитель
забирает с паузой 200~мс. Наконец, организуйте параллельный поиск слова
\texttt{java} в трёх файлах: по одному \texttt{Callable<List<String>>} на файл,
фиксированный пул, сбор результатов через \texttt{Future.get()},
\texttt{shutdown()} в конце.

\subsection*{Результаты занятия}

Должны быть готовы файлы \texttt{.java} по всем четырём заданиям, снимки вывода
и письменные ответы на поставленные в заданиях вопросы.

\kontrolvoprosy

\begin{enumerate}
\item Чем \texttt{ArrayList} отличается от \texttt{LinkedList} и в каких задачах
      выигрывает каждый?
\item Почему \texttt{equals()} и \texttt{hashCode()} переопределяют вместе, если
      объекты будут ключами \texttt{HashMap}?
\item Зачем нужен \texttt{BufferedReader}, если есть \texttt{FileReader}, и чем
      опасен \texttt{FileWriter} без явно заданной кодировки?
\item Чем процесс отличается от потока выполнения? Что у потоков общее, а что
      своё?
\item Почему \texttt{t.run()} не создаёт нового потока, а \texttt{t.start()}~---
      создаёт?
\item Какие две гарантии даёт \texttt{synchronized} и какую из них~---
      \texttt{volatile}?
\item Почему \texttt{wait()} освобождает монитор, а \texttt{sleep()}~--- нет?
\item Чем пул потоков лучше создания \texttt{new Thread()} на каждую задачу?
\end{enumerate}

\testzadaniya

\begin{enumerate}

\vopros{Какое утверждение об интерфейсе \texttt{Map} верно?}
  {наследует \texttt{Collection}}{наследует \texttt{Iterable}}
  {является подтипом \texttt{Set}}
  {не наследует \texttt{Collection}: это отдельная ветвь иерархии}

\vopros{В чём ключевое отличие \texttt{Comparable} от \texttt{Comparator}?}
  {первый задаёт естественный порядок внутри класса, второй~--- стратегия
   сравнения снаружи}
  {первый сравнивает любые типы, второй~--- только один}
  {они различаются только пакетом}
  {первый работает лишь с числами}

\vopros{В пустую \texttt{PriorityQueue} добавили 30, 10 и 20. Что напечатают два
  подряд идущих вызова \texttt{poll()}?}
  {30 и 10}{30 и 20}{10 и 20}{20 и 10}

\vopros{Слова cat, car, dog, deer сгруппировали по первой букве через
  \texttt{Collectors.groupingBy}. Что вернёт группа c?}
  {\texttt{[cat, car]}}{\texttt{[cat, car, dog, deer]}}
  {\texttt{[dog, deer]}}{ошибку компиляции}

\vopros{Какая из операций потока данных терминальная?}
  {\texttt{filter()}}{\texttt{map()}}{\texttt{collect()}}{\texttt{sorted()}}

\vopros{Чем символьные потоки отличаются от байтовых?}
  {работают с \texttt{char} и учитывают кодировку}{всегда быстрее байтовых}
  {байтовые не умеют читать файлы}{работают только с UTF-8}

\vopros{Чему равно \texttt{transient}-поле типа \texttt{String} сразу после
  десериализации объекта?}
  {значению, которое было в момент сохранения}{пустой строке}
  {\texttt{null}}{значению, заданному в конструкторе класса}

\vopros{У объекта \texttt{Thread}, задача которого печатает A, вызвали
  \texttt{run()}, а сразу после этого напечатали B. Каким будет вывод?}
  {AB: вызов \texttt{run()} выполняется в текущем потоке}
  {BA: новый поток выполнится позже}
  {порядок случайный}{возникнет ошибка выполнения}

\vopros{Почему операция \texttt{counter++} не потокобезопасна?}
  {это три действия, между которыми может вмешаться другой поток}
  {операция \texttt{++} применима только к типу \texttt{long}}
  {компилятор оптимизирует эту операцию}
  {она выбрасывает исключение в многопоточной среде}

\vopros{Что произойдёт при вызове \texttt{wait()} вне синхронизированного блока?}
  {поток заснёт и будет ждать \texttt{notify()}}
  {ничего: метод сразу вернёт управление}
  {код не скомпилируется}{будет выброшено \texttt{IllegalMonitorStateException}}

\end{enumerate}
